\section{\bi{Discussion}{讨论}}\label{sec:discussion}
\bi{The two proofs identify a common geometric mechanism with different statistical costs. Magnetic periodization preserves finite positive relative operators exactly; it does not preserve products of arbitrary plane operators. The missing product information is recovered through the full-COM pair frame and the explicit connected-kernel unfolding. This separation explains why local few-body sum-of-squares data can lead to particle-number-uniform torus estimates without a semidefinite program whose dimension grows with flux.}
{两条证明体现了共同几何机制与不同统计代价。磁周期化精确保留有限的正相对算符，但不保持任意平面算符的乘积。乘积中缺少的信息由全COM pair frame及明确的连通核展开补回。这种分离解释了为何局部少体平方和数据能够导出对粒子数一致的环面估计，而不需要维数随磁通增长的半正定规划。}

\bi{The bosonic comparison is not obtained by changing signs in the fermionic proof. Its coherent thermal operators first dominate relative cutoffs and the negative tail as positive few-body coefficients. Their periodizations become integrals over a torus anchor and unrestricted planar spectator offsets. Occupation inequalities then control those integrals by $H$ and $S_4$. Retaining a positive $S_4$ reserve is what permits arbitrarily large occupations. In contrast, Pauli bounds allow the fermionic finite left factors and tail to be paid for by $H$ alone.}
{玻色比较不是在费米证明中改几个符号得到的。其相干thermal算符先在正少体系数层支配相对截断与负尾，周期化后成为环面anchor与不受限平面spectator位移上的积分，再由占据数不等式用$H$及$S_4$控制。保留正$S_4$余量使任意大占据数成为可处理的。相反，Pauli界允许费米有限左因子与尾项都只由$H$支付。}

\bi{The constants are sufficient lower bounds in a fixed physical normalization, not estimates of the actual excitation energies. The bosonic threshold $\ell(\Lambda)\ge20$ is explicit but intentionally coarse. The fermionic threshold remains existential because the high-degree negative-tail control is taken by dominated convergence. A quantitative treatment of that tail would turn this part into an explicit finite-size statement. Improved planar certificates, including the stronger fermionic bounds in~\cite{Gao2026ImprovedGap}, are possible inputs to a further transfer, but their finite matrices and geometric budgets would have to be checked afresh; no such improved torus constant is asserted here.}
{这些常数是在固定物理归一化下的充分下界，不是实际激发能的估计。玻色阈值$\ell(\Lambda)\ge20$显式但有意保守。费米阈值仍是存在性的，因为高次数负尾通过支配收敛控制。若对该尾项作定量处理，就可把这一部分变为显式有限尺寸结论。更好的平面证书，包括文献~\cite{Gao2026ImprovedGap}中更强的费米下界，可能成为进一步迁移的输入，但必须重新检查相应有限矩阵与几何预算；本文不声称这些改进的环面常数已经成立。}

\bi{There are two scope boundaries. First, we do not replace the systole hypothesis by a claim of arbitrary shape uniformity, as discussed in Section~\ref{sec:groundspaces}. Second, this transfer is a quotient construction, not an orbital compression; Appendix~\ref{app:disk} gives an exact disk counterexample and identifies the failing frame estimate. The results also do not prove stability under general perturbations of the torus Hamiltonian. Establishing such stability would require further local-energy and perturbation estimates beyond the finite certificate comparison developed here.}
{适用范围有两条边界。第一，本文不将最短周期条件替换为任意形状的一致性声明，见第~\ref{sec:groundspaces}节。第二，本迁移是商空间构造，不是轨道压缩；附录~\ref{app:disk}给出精确disk反例，并指出失效的frame估计。本文也没有证明环面Hamiltonian在一般扰动下的稳定性。要得到这种稳定性，需要超出这里有限证书比较之外的局部能量与扰动估计。}

\section{\bi{Conclusion}{结论}}
\bi{The bounds below are the mathematical conclusion of the candidate argument, with the formal-verification scope specified in Appendix~\ref{app:reproducibility}. No priority claim is made.}
{下述界是本候选论证的数学结论，形式化验证范围见附录~\ref{app:reproducibility}。本文不作首创性声明。}
\bi{Fixed finite Fock-space certificates can be transferred to full bosonic $V_0$ and fermionic $V_1$ interactions on large-systole flat tori by combining positive periodization, parameter COM integration, and all-COM jet-frame bounds. The separate CCR and CAR comparisons give the gap estimates of Theorems~\ref{thm:boson} and~\ref{thm:fermion}; at $M=2N$ and $M=3N$ they apply above the entire two- and three-dimensional Laughlin ground spaces. The proof depends on large actual cycles and on the specified contact normalization, and does not extend by compression to a hard-cutoff disk.}
{通过结合正周期化、参数COM积分与全COM jet-frame界，可以把固定有限Fock空间证书迁移到大最短周期平坦环面上的完整玻色$V_0$和费米$V_1$相互作用。分别采用CCR和CAR比较，得到定理~\ref{thm:boson}与~\ref{thm:fermion}的谱隙估计；在$M=2N$及$M=3N$时，它们控制整个二维、三维Laughlin基态空间之上的谱。证明依赖实际周期足够大及明确的接触归一化，不能通过压缩延拓到硬截断disk。}
